% !TEX program = xelatex
\documentclass[10pt]{ctexart}
%=====================================================================
% theme-violet.tex —— 讲义精讲 · 紫花主题
% 用法：\documentclass[11pt]{ctexart}
%       \input{assets/theme-violet}
%       \begin{document} ... \end{document}
% 只用 ctexart 自带能力 + xcolor/enumitem/fancyhdr/pifont/amsmath，
% 不依赖 tikz/tcolorbox，保证 MiKTeX 可直接编译。
%=====================================================================

\usepackage[a4paper,left=2.0cm,right=2.0cm,top=2.1cm,bottom=2.3cm,%
            headsep=12pt,footskip=22pt]{geometry}
\usepackage{amsmath,amssymb}
\usepackage{needspace}
\usepackage{xcolor}
\usepackage{array}
\usepackage{enumitem}
\usepackage{pifont}
\usepackage{fancyhdr}
\usepackage[hidelinks]{hyperref}

%---------------------------------------------------------------- 配色
\definecolor{violet}{RGB}{107,78,143}      % 主色：紫
\definecolor{violetdeep}{RGB}{63,44,92}    % 深紫：标题/强调
\definecolor{petal}{RGB}{183,159,208}      % 花瓣：装饰线、花饰
\definecolor{bandviolet}{RGB}{247,243,251} % 浅紫底：档案/提示框
\definecolor{bandcream}{RGB}{252,249,242}  % 米白底：易错/注意框
\definecolor{rulepurple}{RGB}{217,207,230} % 分隔线
\definecolor{inkgray}{RGB}{92,96,110}      % 次要文字

%---------------------------------------------------------------- 字体
% 有 Noto CJK 用 Noto，没有就退回 ctex 默认（Windows 下为宋体/黑体）
\IfFontExistsTF{Noto Serif CJK SC}{\setCJKmainfont{Noto Serif CJK SC}}{}
\IfFontExistsTF{Noto Sans CJK SC}{\setCJKsansfont{Noto Sans CJK SC}}{}

%---------------------------------------------------------------- 版式
\setlist{nosep,leftmargin=1.9em}
\linespread{1.28}
\setlength{\parskip}{3.5pt}
\setlength{\parindent}{0pt}
\setlength{\tabcolsep}{5pt}
\renewcommand{\arraystretch}{1.28}

%---------------------------------------------------------------- 花饰
% \ding{95} = 五瓣小花，\ding{96} = 星形花；都来自 pifont，必定存在
\newcommand{\orn}{\textcolor{petal}{\ding{95}}}          % 行内小标
\newcommand{\ornstar}{\textcolor{petal}{\ding{96}}}      % 星形花
\newcommand{\flowerrule}{%                                % 居中花饰分隔线
  \par\vspace{7pt}%
  \noindent\makebox[\linewidth]{\color{petal}%
    \rule{0.20\linewidth}{0.6pt}\hspace{8pt}{\small\ding{95}}\hspace{4pt}%
    \raisebox{1.1pt}{\tiny\ding{95}}\hspace{4pt}{\small\ding{95}}\hspace{8pt}%
    \rule{0.20\linewidth}{0.6pt}}%
  \par\vspace{7pt}}

%---------------------------------------------------------------- 标题与节头
\newcommand{\doctitle}[3]{%   #1 主标题  #2 副标题  #3 元信息（可用 \\ 分行）
  \begin{center}
    {\color{petal}\small \orn\enspace #2\enspace \ornstar}\\[7pt]
    {\LARGE\bfseries\color{violetdeep} #1}\\[10pt]
    {\color{petal}\rule{0.55\linewidth}{1pt}}\\[4pt]
    {\small\color{inkgray} #3}
  \end{center}\vspace{2pt}}

\newcommand{\hdr}[2]{%   #1 章节号（如 §0）  #2 标题
  \par\vspace{12pt}%
  \noindent\colorbox{violetdeep}{\color{white}\bfseries\small\;#1\;}\hspace{7pt}%
  {\color{violetdeep}\large\bfseries #2}\par\nobreak
  \vspace{3pt}{\color{petal}\hrule height 1.2pt}\vspace{6pt}}

\newcommand{\subhdr}[1]{%
  \par\vspace{6pt}%
  \noindent{\color{violet}\bfseries #1}\par\vspace{1pt}}

%---------------------------------------------------------------- 提示框
\newcommand{\callout}[1]{% 浅紫底：讲义档案、读法说明
  \par\medskip\noindent
  \fcolorbox{rulepurple}{bandviolet}{\parbox{\dimexpr\linewidth-2\fboxsep-2\fboxrule}{\vspace{2pt}#1\vspace{2pt}}}%
  \par\medskip}

\newcommand{\warnbox}[1]{% 米白底：易错、边界条件
  \par\medskip\noindent
  \fcolorbox{rulepurple}{bandcream}{\parbox{\dimexpr\linewidth-2\fboxsep-2\fboxrule}{\vspace{2pt}#1\vspace{2pt}}}%
  \par\medskip}

\newcommand{\gapitem}[2]{% #1 缺口类型（紫） #2 内容
  \par\vspace{2pt}\noindent
  \hangindent=1.4em\hangafter=1
  {\color{petal}\small\ding{96}}\enspace{\bfseries\color{violetdeep}#1}\enspace #2\par}

% 题干块（\qstem）：每道例题／作业的「详解」开头必须先复述题干与出处。
% 起因：2026-10-25 第五章讲义把七道作业的题面集中放在一页、详解另起一页，
% 读者翻到「第 4 题」时看不到条件，等于一道没有题干的题目。
% 规则：详解块内自带题干，读者不翻页就能读懂题在问什么。
\newcommand{\qstem}[1]{%
  \par\medskip\noindent
  \fcolorbox{violet}{bandviolet}{%
    \parbox{\dimexpr\linewidth-2\fboxsep-2\fboxrule}{%
      \vspace{2pt}%
      {\small\bfseries\color{violet}题干}\hspace{0.8em}#1%
      \vspace{2pt}}}%
  \par\medskip}

% v2: 保护题干与开头，长题可跨页，续页补上下文。
\newcommand{\qbreak}{\needspace{6\baselineskip}}
% 独立练习：ID、来源、题干、考点、答案标签（通常 ans:ID）。
\newcommand{\practice}[5]{%
  \qbreak\subhdr{练习 #1}\qstem{（#2）#3}%
  \textbf{考点}：#4\par
  详解见附录第~\pageref{#5}~页。\par\medskip}
% 附录开头：ID、来源、条件摘要；自动产生 ans:ID 标签。
\newcommand{\answerstart}[3]{%
  \qbreak\subhdr{#1\quad 详解}\label{ans:#1}\qstem{（#2）#3}}
\newcommand{\answercontinue}[2]{%
  \needspace{5\baselineskip}\subhdr{#1\quad 续解}%
  {\small\color{inkgray}本页接续：#2}\par\medskip}
% 延迟提供默认值，让旧讲义可在导言区保留自己的 \newcommand{\step}。
\AtBeginDocument{\providecommand{\step}[1]{\needspace{3\baselineskip}\par\medskip\textbf{#1}\par}}

% 前提旁注：只放在「需要理解的那一步」后面，给陈述 + 本题代入，不重证。
\newcommand{\aside}[1]{%
  \par\vspace{1pt}\noindent
  \hangindent=1.6em\hangafter=1
  {\small\color{inkgray}%
    \hspace*{0.6em}{\color{petal}$\triangleright$}\hspace{2pt}#1}%
  \par\vspace{2pt}}

%---------------------------------------------------------------- 表格
% 本机 MiKTeX 的 colortbl 与 array 版本不匹配（\rowcolor 会报错），
% 故表格统一用「上粗紫线 + 表头 + 中细线」的三线式，不依赖 colortbl。
\newcommand{\tabtop}{\par\vspace{3pt}{\color{violet}\hrule height 1.2pt}\vspace{1pt}}
\newcommand{\tabmid}{\par\vspace{1pt}{\color{petal}\hrule height 0.5pt}\vspace{1pt}}
\newcommand{\tabbot}{\par\vspace{3pt}{\color{rulepurple}\hrule height 0.5pt}\vspace{8pt}}
\newcommand{\thead}[1]{{\bfseries\color{violetdeep}#1}}

%---------------------------------------------------------------- 页眉页脚
\pagestyle{fancy}
\fancyhf{}
\renewcommand{\headrulewidth}{0.4pt}
\renewcommand{\headrule}{\hbox to\headwidth{\color{rulepurple}\leaders\hrule height \headrulewidth\hfill}}
\fancyhead[L]{\small\color{inkgray}\songti \doctag}
\fancyhead[R]{\small\color{inkgray}\songti \coursetag}
\fancyfoot[C]{\small\color{inkgray}\orn\enspace{\color{violet}\thepage}\enspace\orn}

% 每篇讲义在导言区设这两个值
\newcommand{\doctag}{}
\newcommand{\coursetag}{}
\newcommand{\settags}[2]{\renewcommand{\doctag}{#1}\renewcommand{\coursetag}{#2}}

%---------------------------------------------------------------- 数学记号
\newcommand{\vect}[1]{\boldsymbol{#1}}
\newcommand{\dd}{\,\mathrm{d}}
\newcommand{\ddx}{\frac{\mathrm{d}}{\mathrm{d}x}}
\addtolength{\jot}{2pt}



\usepackage{graphicx}
\setlength{\headheight}{20pt}
\settags{实验 5 · 速查卡}{电工电子技术}
\pagestyle{empty}
\begin{document}

\begin{center}
{\LARGE\bfseries\color{violetdeep} 实验 5 · 单管放大电路 速查卡}\\[4pt]
{\small\color{inkgray} 2026-11-05（周四）15:30--17:45 \quad 综合一号教学楼 317 \quad 2502 班 \quad 10 分}
\end{center}

\vspace{-2pt}
{\color{petal}\hrule height 1pt}
\vspace{4pt}

\textbf{\color{violetdeep}一、出门前}：教材《电工学实验教程》第 2 版、\textbf{预习报告}（没有不许做实验）、报告纸、铅笔直尺、计算器。

\textbf{\color{violetdeep}二、四组偏置的静态工作点（算例：$U_{CC}=12$ V、$\beta=100$、$U_{BE}=0.7$ V、$S_1$ 断开）}

\begin{center}
\small
\begin{tabular}{@{}p{1.2cm}p{1.3cm}p{1.3cm}p{1.2cm}p{1.3cm}p{1.4cm}p{1.4cm}p{2.3cm}p{2.2cm}@{}}
\hline
\thead{组合} & \thead{$R_{B1}$} & \thead{$R_{B2}$} & \thead{$V_{BB}$} & \thead{$I_B$} & \thead{$I_C$} & \thead{$U_{CE}$} & \thead{区域} & \thead{现场判断} \\
\hline
1 & 30 k & 39 k & 6.78 V & 44 $\mu$A & 4.40 mA & $-6.49$ V & 饱和 & $U_{CE}\approx0.3$ V \\
2 & 30 k & 12 k & 3.43 V & 21 $\mu$A & 2.10 mA & 3.17 V & 放大·偏低 & 下半周先削 \\
3 & 150 k & 39 k & 2.48 V & 11.7 $\mu$A & 1.17 mA & 7.10 V & \textbf{放大·最佳} & 用这组做动态 \\
4 & 150 k & 12 k & 0.89 V & 1.43 $\mu$A & 0.143 mA & 11.40 V & 放大·近截止 & 上半周先削 \\
\hline
\end{tabular}
\end{center}

\textbf{\color{violetdeep}三、三个公式（第 3 组：$r_{be}=200+100\times26/1.17\approx2.43$ k$\Omega$）}

\[
A_u=-\frac{\beta R_L'}{r_{be}}=-61.8\ (R_L=3\ \mathrm{k}\Omega),\qquad
r_i=R_B\parallel r_{be}=2.25\ \mathrm{k}\Omega,\qquad
r_o\approx R_C=3\ \mathrm{k}\Omega .
\]

\textbf{\color{violetdeep}四、$S_1$ 一开一合（$R_{E2}=200\,\Omega$）}

\begin{center}
\small
\begin{tabular}{@{}p{2.4cm}p{4.3cm}p{4.3cm}p{4.5cm}@{}}
\hline
\thead{状态} & \thead{$A_u$（$R_L=3$ k）} & \thead{$r_i$} & \thead{其他} \\
\hline
$S_1$ 闭合（无反馈） & $-61.8$ & 2.25 k$\Omega$ & 失真大、频带窄 \\
$S_1$ 断开（有反馈） & $-6.63$ & 13.1 k$\Omega$ & 失真小、频带宽 \\
\hline
\end{tabular}
\end{center}

有反馈时 $A_u=-\dfrac{\beta R_L'}{r_{be}+(1+\beta)R_{E2}}$。

\textbf{\color{violetdeep}五、报告六栏（英文照抄）}：\textit{I. Experimental Objectives}｜\textit{II. Experimental Apparatus}｜\textit{III. Experimental Principles}｜\textit{IV. Experimental Procedure}｜\textit{V. Results and Analysis}｜\textit{VI. Discussion and Conclusions}

\textbf{\color{violetdeep}六、现场三条纪律}

\begin{enumerate}
\item \textbf{静态先短路输入}：$u_i=0$、输入对地短接；所有表共地；电流不拆线，测 $U_B$、$U_C$、$U_E$ 再算。
\item \textbf{动态先看波形}：示波器确认不削顶再读 $U_o$；$u_i$ 与 $u_o$ 反相 $180^\circ$；换负载、串电阻时\textbf{输入幅度不变}。
\item \textbf{数字如实写}：$f_H$ 在 20 kHz 内测不到就写“未能测出”；饱和组写明“$\beta I_B$ 不成立”，不编数。
\end{enumerate}

\vspace{2pt}
{\small\color{inkgray} 讲义全文（20 页，含六张数据表、术语表、预习与讨论题详解）：\texttt{实验5-单管放大电路-实验讲义-20261104.pdf}；\textbf{电路图}（图 5.2.4 原图 + 端子表 + 报告里怎么画）：\texttt{实验5-电路图页-20261105.pdf}，本卡不重复画图。}

\end{document}
